\section{Discussion}
\label{sec:discussion}

\subsection{What failed and what remains possible}
\label{sec:discussion-failure}

The evidence chain did not fail at a single common point.
Controlled brain predictivity passed in both recorded-response assays \evd{E002}\evd{E006}.
Compression left observed alignment headroom, although language-model quality remained a competing explanation for that headroom \evd{E003}\evd{E015}.
The recorded-target branch then lost support at the intervention gate: none of the valid inference units established a participant-general advantage from correct stimulus--response pairing \evd{E008}\evd{E011}\evd{E013}\evd{E017}.
In the predicted-target branch, the exact target failed the frozen practical measurability rule, while target specificity against the frozen twin remained unresolved \evd{E030}.
Target uptake and retained-student movement passed in the saved WikiText continuity assay at acceptable language quality \evd{E016}\evd{E025}.
Fresh Tuckute incremental retention remained unresolved \evd{E030}.
The saved nonbrain comparator was non-identifying \evd{E026}.
Participant-general biological transfer was not supported \evd{E025}\evd{E030}.
No tested route subsequently established a brain-specific neural-target utility claim \evd{E009}.
This pattern separates a measurable signal and a successful manipulation from an identified biological training benefit.

For recorded neural targets, the participant-first result weakens the interpretation that positive trends for an averaged response generalize to people.
Response averaging remains useful for measuring a shared, higher-signal component, but it changes the estimand and does not create participant replicates \evd{E005}\evd{E010}\evd{E014}.
The intervention variants narrow the failure further without making it objective- or capacity-invariant.
Increasing adapter rank at a fixed loss weight did not induce greater effective movement, so it tests added adapter capacity under the same optimization pressure rather than a stronger manipulation.
The contrastive result is limited to the tested five-\ac{roi} objective, while the voxelwise and full-parameter results close only their tested naturalistic data, loss, and optimization regimes.
They weaken those specific rescue explanations, not the broader possibility that another recorded target or induction regime could work \evd{E008}\evd{E011}\evd{E013}\evd{E017}.

For predicted neural targets, the exact target did not meet the frozen practical criterion for unique linear predictivity of the intended recorded-response substrate, and its row-specific advantage over the frozen twin remained unresolved \evd{E030}.
In the saved WikiText continuity assay, target uptake and retained-student movement passed after the temporary head was removed, and the direct comparisons preserved language-model quality, making a simple implementation failure less plausible \evd{E016}\evd{E025}.
Fresh Tuckute incremental retention nevertheless remained unresolved \evd{E030}.
The saved nonbrain target also differed in baseline predictability, remaining headroom, covariance geometry, nuisance predictability, and induced movement, so the larger target-learning gain could not be attributed to neural target content \evd{E026}.
The near-zero direct participant endpoint then supplied no participant-general transfer evidence \evd{E025}.
The first defensible failure is therefore before biological transfer, but neither the measurability result nor the comparator audit explains the later endpoint causally.

Several explanations are consequently weakened, whereas others remain live.
The nuisance and untrained-network controls weaken the claim that no controlled neural signal was available to the declared linear readouts \evd{E002}\evd{E006}.
The target uptake, retained-student movement, and quality checks weaken no-op training and gross language-quality damage as explanations for the saved predicted-target result \evd{E016}\evd{E025}.
They do not distinguish insufficient target information from unfavorable target geometry, a mismatch between slow or averaged \ac{fmri} targets and the representations being trained, inadequate temporal precision, poorly directed optimization pressure, limited model or data scale, information accessible only through a nonlinear map, or heterogeneous participant effects.
These are hypotheses suggested by where support was lost, not findings about why it was lost.
The conclusion is therefore a scoped localization over the tested targets, transforms, objectives, models, participants, controls, linear assays, and endpoints, not a universal null for neural supervision.

\subsection{Compatibility with positive prior studies}
\label{sec:discussion-prior}

The local negative result is compatible with positive neural-guided studies because the interventions and estimands are not interchangeable.
As Section~\ref{sec:prior-gap} details, the literature spans speech and text models, slow \ac{fmri} and temporally resolved recordings, participant-specific and pooled supervision, adaptation and compression, and biological and behavioral endpoints.
A high-temporal-resolution target can constrain token- or word-scale representations differently from a sentence-level or hemodynamically blurred target.
Personalized training can exploit stable within-person structure that an across-participant claim must treat as heterogeneity.
Likewise, a larger retained encoder, a more deeply sampled participant, or a larger training budget can make a target learnable under conditions that need not hold for a fixed smaller student.
These differences change the intervention being tested rather than merely changing its implementation details.

The claim licensed by a positive contrast also depends on its comparison and analysis.
A gain over a pretrained model answers a different question from an incremental gain over ordinary \ac{kd} at comparable language-model quality.
A gain over a valid, marginal-preserving same-target permutation supports correct stimulus--target pairing but does not distinguish neural target content from other correctly paired, learnable stimulus-derived targets.
The defective saved block permutation does not meet that condition and therefore cannot support clean pairing attribution \evd{E016}.
Repeated training seeds establish procedural stability, whereas participant-level uncertainty is required for a claim about people; folds and neural measurement units answer still different sampling questions.
Endpoint choice matters in the same way: improvement on the optimized neural target is target uptake, improvement on newly recorded responses is biological transfer, and improvement on a task is utility.
A positive result can be valid for any one of these estimands without satisfying all later gates.

The present study therefore does not refute positive results obtained with different modalities, participant designs, model budgets, baselines, or endpoints.
It adds a narrower compatibility judgment: controlled brain predictivity can coexist with an unsupported training advantage, and successful target uptake can coexist with a non-identifying comparator and absent participant-general transfer.
Its contribution is to keep those claims separate and to show which controls and inference units are needed when moving from one to the next.
Whether a positive prior result would survive every additional gate is an open empirical question, not a conclusion licensed by the present experiments.

\subsection{Design lessons from the observed failures}
\label{sec:discussion-design}

The comparator failure makes target matching a design problem that must be solved before attribution to neural target content.
A nonbrain target should match the neural target's accessible stimulus information as closely as the scientific question permits, along with dimensional scale, covariance geometry, nuisance predictability, baseline predictability, remaining headroom, held-out learnability, and the optimization pressure applied to the student.
Initial auxiliary losses and gradient scales should be retained so that pressure can be audited rather than inferred after training.
Equal output width or a common training schedule is insufficient when one target is already easy for the \ac{kd} student or induces substantially different movement.
This recommendation follows from the observed non-comparability; it does not imply that one fixed matching recipe will identify neural target content in every setting \evd{E026}.

The exact-target failure also changes the order of work for predicted neural targets.
Before student training, the exact transformed target should be tested on the intended recorded-response substrate using held-out data, the planned nuisance model and readout, an appropriate frozen target control, participant-level inference when the claim concerns people, and a predeclared threshold large enough to matter for the downstream study.
The precheck should require both practically measurable unique signal beyond nuisance and adjudicated target specificity relative to that control.
In the present assay, the unique increment failed the frozen practical rule and aligned-versus-twin specificity remained unresolved \evd{E030}.
Failure at either part should trigger a target, transform, modality, control, or assay redesign before expensive optimization.
Passing both would not guarantee transfer, but it would establish under that assay a specific measurable signal that the later biological transfer test could in principle preserve.
The observed failure supports this ordering for diagnostic efficiency, not a universal claim that linear measurability is necessary for every possible nonlinear intervention.

For predicted neural targets, the decision rule is conjunctive: exact-target measurability must pass; target uptake and retained-student movement must then pass at acceptable language quality; comparator validity must be established, including matched trained-model learnability and movement; and biological transfer must be supported on an independent recorded endpoint at participant level \evd{E016}\evd{E025}\evd{E026}\evd{E030}.
After target uptake and retained-student movement pass, evaluation must move away from the optimized target.
The direct target-trained-versus-\ac{kd} contrast should be measured on independently recorded responses, with technical seeds and stimulus folds aggregated within each participant before population inference.
A valid matched comparator is additionally required to attribute any transfer to neural target content rather than generic auxiliary supervision.
Reporting target uptake, retained-student movement, language quality, and participant-level transfer as separate outcomes makes the failure location visible: a manipulation can be reliable even when its biological endpoint is not.
Only after every prerequisite passes does a downstream utility test estimate the payoff of an identified biological intervention \evd{E025}.

The gate location should determine whether a study continues, redesigns, or stops.
When target uptake and retained-student movement pass but participant-level transfer remains unresolved, strengthening or repeating the transfer test is justified only if added participants, a better independent endpoint, or a more informative assay could change the decision.
When a practically meaningful effect is excluded or an upstream gate fails, the warranted choices are to stop or redesign the target, control, or assay rather than accumulate more target-fitting evidence.
Utility testing should wait in either case.
This decision rule avoids spending downstream compute to interpret an upstream failure and avoids treating every nonpass as the same kind of null.
It is a recommendation derived from the diagnostic pattern in this study, not an experimentally proven universal recipe for brain-guided training.
